\subsection{The signed CM character on the full GT group}

\begin{result}{E.191}{Proved}
{One CM exponent forces the two historical frame parities to agree}
The historical continuous character
\(\Delta:\widehat{\mathrm{GT}}_0\to\mathbb F_2\),
with the literal signed pure frames of E.95, satisfies
\[
                       \Delta|_{\widehat{\mathrm{GT}}}=0 .
\tag{E.191.1}
\]
Thus the vanishing clause of historical GT A.7 and
the signed CM comparison in its A.3 hold.  The proof
uses the actual full completions and the single
comparison of E.190.  It does not replace that
comparison by common division fields or separate
conjugators modulo three and four.
\end{result}

\begin{proof}
By E.181 and the arithmetic vanishing and character
law proved in E.95, it suffices to consider
\(F=(1,f)\in\widehat{\mathrm{GT}}\) with \(\rho(F)=0\).
Keep E.190's \(\alpha,\beta,A,\Phi,\varphi\).  Write
\[
 \begin{gathered}
 U=\begin{pmatrix}1&1\\0&1\end{pmatrix},\qquad
 L=\begin{pmatrix}1&0\\-1&1\end{pmatrix},\\
 S=ULU,\qquad R_+=UL,\qquad R_0=-R_+,\qquad
 T=q_{12}(A)^{-1},\quad W=q_{12}(\beta)^{-1}.
 \end{gathered}
\tag{E.191.2}
\]
The notation \(R_+\) denotes both the actual modular
element \(uw\) and its displayed homology matrix.
The signed pure evaluation is
\(x\mapsto U^2,\ y\mapsto L^2\).
E.95.6--E.95.7, now with zero shear, give
\[
 \begin{gathered}
 q_{12}(f)=I,\qquad T=W=I\pmod4,\\
 T\bmod3\in C_{\operatorname{SL}_2(\mathbb F_3)}(R_0),
 \qquad W\bmod3\in\langle S\rangle,\\
 \delta(F)=\epsilon_3(T)+\epsilon_2(W).
 \end{gathered}
\tag{E.191.3}
\]
Here \(\epsilon_3\) is the scalar sign on
\(\ker(R_0-I)\), and \(\epsilon_2(S^i)=i\bmod2\).
These are exactly the signed historical bits;
the equality \(V=I\) is a consequence of normalization,
not a changed definition.

\emph{The target homology and its exact frame.}
Use the ordinary presentations
\[
 \Gamma_{1,2}=B_4/\langle(abc)^4\rangle,\qquad
 \Gamma_{1,1}=B_3/\langle(uw)^6\rangle,
 \qquad(a,b,c)\longmapsto(u,w,u).
\tag{E.191.4}
\]
The point-pushing kernel is free on
\[
       \xi_1=ca^{-1},\qquad
       \xi_2=aba c^{-1}b^{-1}a^{-1}.
\tag{E.191.5}
\]
For precision, the exact conjugation calculation is
\[
 \begin{array}{c|cc}
       &\xi_1&\xi_2\\ \hline
 a&\xi_1&\xi_2\xi_1\\
 b&\xi_1\xi_2^{-1}&\xi_2\\
 c&\xi_1&\xi_1\xi_2 .
 \end{array}
\tag{E.191.6}
\]
This can also be obtained by a presentation
calculation: introduce \(x=ca^{-1}\) and
\(y=abx^{-1}b^{-1}a^{-1}\) in the braid
presentation.  The relations \(ac=ca\) and
\(bcb=cbc\), together with \(aba=bab\), are
equivalent to the four conjugation relations
for \(a,b\) in E.191.6.  Conversely, those
relations with \(c=xa\) give the three
four-strand braid relations.  This is the
semidirect presentation \(B_4=F(x,y)\rtimes B_3\);
the \(c\) row is conjugation by \(xa\).
The central word \((abc)^4\) maps to the
infinite-order word \((uw)^6\), so quotienting
by the indicated central subgroups retains
exactly that free kernel.  E.182 retains its
full topology after completion.
Thus \(a,c\) act by \(U\), and \(b\) by \(L\),
on \(K^{\rm ab}=\widehat{\mathbb Z}^2\).

Put \(e=f(ba^2b,c^2)\).  The actual standard
four-strand formulas are
\[
 \Phi(a)=a,\qquad
 \Phi(b)=f(a^2,b^2)^{-1}b f(a^2,b^2),\qquad
 \Phi(c)=e^{-1}ce .
\tag{E.191.7}
\]
Since \(ba^2b=(ab)^3a^{-2}\), the removed
factor \(a^{-2}\) commutes with both remaining
arguments.  A commutator proword has zero
exponent sums, so
\(e=f((ab)^3,c^2)\) and \([e,a]=1\).
After applying \(q\) those two arguments commute,
so \(e\in K\).

Let \(N\) be the induced matrix of \(\Phi\) on
\(K^{\rm ab}\), and let \(E=[e]\).
The preceding equalities give
\[
       (U-I)E=0,\qquad
       Ne_1=e_1+(U-I)E=e_1,\qquad NU=UN .
\tag{E.191.8}
\]
Writing a two-by-two matrix shows that \(N=U^t\)
for some \(t\in\widehat{\mathbb Z}\).
Covariance of the second formula of E.191.7 gives
\[
 NLN^{-1}
       =q_{12}(f)^{-1}Lq_{12}(f)=L\pmod{12}.
\tag{E.191.9}
\]
A matrix commuting with both \(U\) and \(L\)
over \(\mathbb Z/12\) is scalar: the first
commutation sets its lower-left entry to zero
and its diagonal entries equal, and the second
sets its upper-right entry to zero.
Since \(N\) fixes \(e_1\), this scalar is one.
Therefore \(N=I\pmod{12}\).
The lift \(\widetilde A=A(a^2,b^2)\) in E.190
acts on \(K^{\rm ab}\) by the actual signed
matrix \(q_{12}(A)\).  It follows that
\[
              (\Phi^A)_{K^{\rm ab}}=T\pmod{12}.
\tag{E.191.10}
\]
This identifies the target frame including its
modulo-four sign.

\emph{The source homology and its two residues.}
Let \(K_0=\mathcal B'\).  Its full metabelian
quotient has the exact form
\[
 \begin{gathered}
 \mathcal B/[K_0,K_0]
      =\widehat{\mathbb Z}^2\rtimes C_6,\qquad
 R=\begin{pmatrix}0&-1\\1&1\end{pmatrix},\\
 \tau=(0,R),\quad s=(p,R^3),\quad
 r=(p,R^4),\qquad p=e_1 .
 \end{gathered}
\tag{E.191.11}
\]
Here is a check of the splitting and the integral
module, rather than an assumed rational model.
In the index-six Schreier presentation put
\(p_i=\tau^i(s\tau^{-3})\tau^{-i}\) and
\(z=\tau^6\).  The two factor relations give,
after abelianization,
\[
 p_i+p_{i+3}+z=0,\qquad
 p_i+p_{i+4}+p_{i+8}+2z=0 .
\tag{E.191.12}
\]
Indices shifted by six have the same abelian
class.  The first relations express \(p_3,p_4,p_5\)
using \(p_0,p_1,p_2,z\).  The second ones for
successive \(i\) give \(z=0\) and
\(p_2=p_1-p_0\).  The displayed semidirect
product satisfies both factor relations and
has independent \(p_0,p_1\), proving that there
is no additional relation.
The actual finite-index topology then gives
E.191.11 over \(\widehat{\mathbb Z}\).

Let \(M\) be the matrix of \(\alpha\) on
\(K_0^{\rm ab}\).  Since \(\alpha(\tau)=\tau\),
\(MR=RM\).  Write
\(\beta=(b,R^j)\), \(A=(a,R^\ell)\).
Purity makes \(j,\ell\) even, and
\((\beta^{-1})^{\rm ab}=(\rho,\rho)=0\)
in the original pure two-generator group
makes \(j=0\).  The literal frame equations
now give
\[
          Mp=2b+p=(I-R^4)a+R^\ell p .
\tag{E.191.13}
\]
Modulo three \(R^4=-R\), and \(R+I\)
has rank one, square zero, and equal image
and kernel \(L_-\).  The second expression
of E.191.13 shows that \(M\) induces one
on the quotient by \(L_-\), because
\(\ell\) is even.  Commutation with \(R+I\)
then gives one on the line itself:
\[
                         M|_{L_-}=1\pmod3 .
\tag{E.191.14}
\]

For the integral modulo-four sign, identify
the projective action of
\(\operatorname{PSL}_2(\mathbb F_3)=A_4\)
with the affine action on \(\mathbb F_2^2\)
by the literal assignment
\[
 \infty\mapsto0,\quad 0\mapsto p,\quad
 -1\mapsto Rp,\quad 1\mapsto p+Rp .
\tag{E.191.15}
\]
Direct application of the fractional linear
maps gives
\[
             S=(p,I),\qquad R_0=(p,R^4),
                       \qquad U^{-1}=(0,R).
\tag{E.191.16}
\]
Consequently E.191.11 modulo two is exactly
this projective modular representation.
As \(j=0\), the projective image of
\(W=q_{12}(\beta)^{-1}\) is translation by
\(-b=b\).  Its membership in \(\langle S\rangle\)
therefore says
\[
                        b=\epsilon_2(W)p\pmod2 .
\tag{E.191.17}
\]
Substituting in the first expression of
E.191.13 and using commutation with \(R\)
on the basis \(p,Rp\) proves
\[
                         M=(-1)^{\epsilon_2(W)}I\pmod4 .
\tag{E.191.18}
\]
This step retains the full signed source
frame; projectivization was used only to
compute \(b\) modulo two.

\emph{The common exponent.}
Write \(q(h)=R_+^k\) for the single element
in E.190.4; thus \(k\in\mathbb Z/6\).
The actual discrete map \(\varphi\) gives
an integral isomorphism
\[
 C:K_0^{\rm ab}\xrightarrow{\sim}K^{\rm ab},
           \qquad CR=R_+^\varepsilon C,
                       \quad\varepsilon=\pm1 .
\tag{E.191.19}
\]
Inner elements of \(K\) act trivially on
its abelianization.  Applying it to the
one covariance equation gives
\[
                          TC=R_+^k CM\pmod{12}.
\tag{E.191.20}
\]
At three, \(C\) sends \(L_-\) onto
\(\ker(R_0-I)\); either value of
\(\varepsilon\) has the same minus-one
eigenline.  By E.191.14 the right side
acts on that line with scalar \((-1)^k\).
Hence \(\epsilon_3(T)=k\bmod2\).

At four, E.191.3 and E.191.18 turn
E.191.20 into
\[
                      R_+^k=(-1)^{\epsilon_2(W)}I .
\tag{E.191.21}
\]
The matrix \(R_+\) has order six modulo
four and \(R_+^3=-I\).  Therefore
\(k\equiv3\epsilon_2(W)\pmod6\), so
\(\epsilon_2(W)=k\bmod2\).
The two bits agree, proving \(\delta(F)=0\)
on the normalized subgroup.
E.181 writes any full GT element as
such an element times an arithmetic one.
E.95 supplies both multiplicativity and
arithmetic vanishing of the same \(\Delta\).
This proves E.191.1 for every multiplier.

For reproducibility,
\path{code/check_gt_cm_covariance.py} independently
checks all six identities in E.191.6 by exact
positive braid defining-relation paths after
explicit Garside denominator clearing.  It records
306 local paths with 295 relation steps, the twelve
projective-affine label checks, both compatible
normalized exponents, and four exact quartic
coordinate identities.  Its canonical certificate
hash is
\[
 \begin{gathered}
 \texttt{d7801fb07aff31f355d618771d723a757}\\
 \texttt{e477fc1d773210ae4cacab148fe0d9a}.
 \end{gathered}
\]
These finite checks verify the displayed algebra;
the full source covariance is the separate
geometric proof of E.190.
\end{proof}
